\subsection{直和与乘积的对偶}

对每个 \(i \in I\)，设 \((\mathbf{E}_{i}, \mathbf{F}_{i})\) 是一对向量空间，由双线性型 \(B_{i}\) 置于对偶之中。令 \(E = \prod_{i \in I} E_{i}\)，\(F = \oplus_{i \in I} F_{i}\)，并把每个 \(F_{i}\) 等同于 F 的一个子空间。我们借助双线性型

\[
\mathrm{B} (x, y) = \sum_ {i \in \mathrm{I}} \mathrm{B} _ {i} \left(x _ {i}, y _ {i}\right) \quad \text { for } \quad x = \left(x _ {i}\right) \quad \text { and } \quad y = \left(y _ {i}\right)\tag{8}
\]

将 E 与 F 置于对偶之中（族 \((\mathbf{B}_i(x_i, y_i))_{i \in I}\) 具有有限支集）。

我们回忆（II，第 50 页，命题 8）：弱拓扑 \(\sigma(\mathrm{E},\mathrm{F})\) 是诸弱拓扑 \(\sigma(\mathrm{E}_i,\mathrm{F}_i)\) 的乘积拓扑。

引理 3. --- (i) 对每个 \(i \in \mathbf{I}\)，设 \(\mathfrak{S}_i\) 是 \(\mathrm{F}_i\) 的子集的一个族，其中的子集对于 \(\sigma(\mathrm{F}_i, \mathrm{E}_i)\) 有界；令 \(\mathfrak{S} = \bigcup_{i \in \mathbf{I}} \mathfrak{S}_i\)。则 \(\mathfrak{S}\)-拓扑在 \(\mathbf{E}\) 上就是诸 \(\mathfrak{S}_i\)-拓扑在 \(\mathrm{E}_i\) 上的乘积拓扑。

\begin{enumerate}
\def\labelenumi{(\roman{enumi})}
\setcounter{enumi}{1}
\tightlist
\item
  对每个 \(i \in \mathbf{I}\)，设 \(\mathfrak{J}_i\) 是空间 \(\mathrm{E}_i\)（赋以弱拓扑 \(\sigma(\mathrm{E}_i, \mathrm{F}_i)\)）上的一个适配有界结构，且它们都不等于 \(\{\varnothing\}\)。设 \(\mathfrak{J}\) 是这样的子集 \(\mathbf{A}\)（属于 \(\mathrm{E} = \prod_{i \in \mathbf{I}} \mathrm{E}_i\)）组成的族：\(\mathrm{pr}_i(\mathbf{A}) \in \mathfrak{J}_i\)（对所有 \(i \in \mathbf{I}\)）。则 \(\mathfrak{J}\)-拓扑在 \(\mathbf{F}\) 上就是诸 \(\mathfrak{J}_i\)-拓扑在 \(\mathrm{F}_i\) 上的直和拓扑。
\end{enumerate}

设 \(\mathcal{T}\) 为诸 \(\mathfrak{S}_i\)-拓扑的乘积拓扑。形如

\[
\mathbf {A} = \prod_ {i \in \mathbf {J}} \mathbf {A} _ {i} ^ {\circ} \times \prod_ {i \in \mathbf {I} - \mathbf {J}} \mathbf {E} _ {i}
\]

（其中 \(\mathbf{J} \subset \mathbf{I}\) 有限，且 \(\mathbf{A}_i \in \mathfrak{S}_i\)（对所有 \(i \in \mathbf{J}\)））的集构成 \(\mathcal{T}\) 的一个基本零邻域系。我们有 \(\mathbf{A} = (\bigcup_{i \in \mathbf{J}} \mathbf{A}_i)^\circ\)，故 \(\mathcal{T}\) 与 \(\mathfrak{S}\)-拓扑相同。这就证明了 (i)。

我们把 \(\mathfrak{J}\)-拓扑赋给 \(F\)，把 \(\mathfrak{J}_i\)-拓扑赋给每个 \(F_i\)。对每个子集 \(A\)（\(E\) 中的），我们有 \(F_i \cap A^\circ = \text{pr}_i(A)^\circ\)，故从 \(F_i\) 到 \(F\) 中的单射连续。设 \(q\) 是 \(F\) 上的半范数；我们假定 \(q_i\)（即 \(q\) 在 \(F_i\) 上的限制）对所有 \(i \in I\) 连续。于是可以找到非空子集 \(A_i \in \mathfrak{J}_i\)，使得我们有

\[
q _ {i} \left(y _ {i}\right) \leqslant \sup _ {x _ {i} \in \mathrm{A} _ {i}} \left| \mathrm{B} _ {i} \left(x _ {i}, y _ {i}\right) \right| \quad \left(y _ {i} \in \mathrm{F} _ {i}\right).\tag{9}
\]

令 \(\mathrm{A} = \prod_{i\in \mathrm{I}}\mathrm{A}_i\)；则 \(\mathrm{A}\in \mathfrak{J}\)。对 \(y = (y_i)_{i\in \mathrm{I}}\)（属于 \(\mathrm{F}\)），我们有

\[
q (y) \leqslant \sum_ {i \in \mathbf {I}} q _ {i} (y _ {i}) \leqslant \sum_ {i \in \mathbf {I}} \sup _ {x _ {i} \in \mathrm{A} _ {i}} | \mathrm{B} _ {i} (x _ {i}, y _ {i}) | = \sup _ {x \in \mathrm{A}} | \mathrm{B} (x, y) |,
\]

其中最后的等式由 (8) 得出，因为族 \((y_{i})_{i\in\mathbf{I}}\) 具有有限支集，且诸 \(A_{i}\) 非空并可设为均衡的（GT，IV，§ 5，第 7 目，命题 12 的推论 2）。这一不等式表明 q 在 F 上连续，由此得 (ii)。

命题 12. --- 拓扑 \(\beta(\mathrm{F},\mathrm{E})\) 是诸拓扑 \(\beta(F_i,E_i)\) 的直和拓扑。拓扑 \(\beta(\mathrm{E},\mathrm{F})\) 是诸拓扑 \(\beta(\mathrm{E}_i,\mathrm{F}_i)\) 的乘积拓扑。

我们将应用引理 3，对 \(\mathfrak{S}_i\) 取 \(F_i\) 中对于 \(\sigma(F_i, E_i)\) 有界的所有子集组成的族，对 \(\mathfrak{J}_i\) 取 \(E_i\) 中对于 \(\sigma(E_i, F_i)\) 有界的所有子集组成的族。

由 III 第 4 页的推论 2，\(\mathfrak{J}\) 是 \(E_{i}\) 中对于诸 \(\sigma(\mathrm{E}_{i}, \mathrm{F}_{i})\) 的乘积拓扑有界的所有子集组成的族，而该乘积拓扑与 \(\sigma(\mathrm{E}, \mathrm{F})\) 相同。于是我们关于 \(\beta(\mathrm{F}, \mathrm{E})\) 的断言得证。

我们对 \(\mathrm{F} = \bigoplus_{i\in \mathrm{I}}\mathrm{F}_i\) 赋以拓扑 \(\mathcal{T}\)，它是诸拓扑 \(\sigma (\mathrm{F}_i,\mathrm{E}_i)\) 的直和拓扑。则 \(\mathrm{F}\) 的对偶由线性型 \(y\mapsto \mathrm{B}(x,y)\) 组成，其中 \(x\) 取遍 \(\mathrm{E}\)（II，第 30 页，命题 6）。由 IV 第 1 页的命题 1，拓扑 \(\mathcal{T}\) 与 \(\sigma (\mathrm{F},\mathrm{E})\) 具有相同的有界集。先设诸拓扑 \(\sigma (\mathrm{F}_i,\mathrm{E}_i)\) 是 Hausdorff 的。由 III 第 5 页的命题 5，这些集包含在形如 \(\sum_{i\in \mathrm{J}}\mathrm{B}_i\) 的子集中，其中 \(J\subset I\) 有限，且 \(B_{i}\) 在 \(F_{i}\) 中（对于 \(\sigma (\mathrm{F}_i,\mathrm{E}_i)\)）有界（对所有 \(i\in J\)）。由于 \(\sum_{i\in J}\mathrm{B}_{i}\) 包含在 \(\bigcup_{i\in J}n\mathrm{B}_{i}\) 的凸包中，其中 \(n = \operatorname {Card}(J)\)，我们可以应用引理 3 来证明此情形中关于 \(\beta (\mathrm{E},\mathrm{F})\) 的断言。

对一般情形，设 \(\mathbf{N}_i\) 是 \(\sigma (\mathrm{F}_i,\mathrm{E}_i)\) 的所有零邻域的交，并令 \(\mathbf{N} = \sum_{i\in \mathbf{I}}\mathbf{N}_i\)；则 \(\mathbf{F} / \mathbf{N}\) 是诸 \(\mathrm{F}_i / \mathrm{N}_i\) 的拓扑直和（II，第 31 页，命题 8）；由此我们推出，\(\mathbf{F}\) 中每个对于 \(\mathcal{T}\) 有界的子集都包含在形如 \(\mathbf{N} + \sum_{i\in \mathbf{J}}\mathbf{B}_i\) 的集中，其中 \(\mathrm{J}\subset \mathrm{I}\) 有限，且 \(\mathbf{B}_i\) 在 \(\mathbf{F}_i\) 中有界（对所有 \(i\in \mathbf{J}\)）（III，第 2 页，注 3）；由于该集在 \(\mathbf{E}\) 中的极集与 \(\sum_{i\in \mathbf{J}}\mathbf{B}_i\) 的极集相同，结论同上得出。

命题 13. --- Mackey 拓扑 \(\tau(\mathrm{F},\mathrm{E})\) 是诸 Mackey 拓扑 \(\tau(\mathrm{F}_i,\mathrm{E}_i)\) 的直和拓扑。拓扑 \(\tau(\mathrm{E},\mathrm{F})\) 是诸拓扑 \(\tau(\mathrm{E}_i,\mathrm{F}_i)\) 的乘积拓扑。

关于 \(\tau(\mathrm{F},\mathrm{E})\) 的断言由引理 3 (ii) 及下列性质得出：\(F^{*} = \prod_{i\in I}F_{i}^{*}\) 的闭、凸且均衡的子集为对于 \(\sigma(F^{*},F)\) 紧，必要且充分的是它在每个 \(\mathrm{F}_i^*\) 上的投影对于 \(\sigma(\mathrm{F}_i^*, \mathrm{F}_i)\) 紧。

为证明关于 \(\tau(\mathrm{E},\mathrm{F})\) 的断言，先设诸拓扑 \(\sigma(\mathrm{F}_i,\mathrm{E}_i)\) 是 Hausdorff 的，只需（引理 3 (i)）证明：F 中对于 \(\sigma(\mathrm{F},\mathrm{E})\) 凸、均衡且紧的每个子集 A 都包含在形如 \(\sum_{i\in\mathbf{J}}\mathrm{A}_i\) 的集中，其中

\(J \subset I\) 有限，且 \(A_{i}\) 对于 \(\sigma(F_{i}, E_{i})\) 凸、均衡且紧。但这样的子集对于 \(\sigma(F, E)\) 有界。由命题 12 的证明，存在有限子集 \(J\)（\(I\) 中的）使得 \(A \subset \sum_{i \in J} F_i\)，而只需对 \(A_i\) 取 \(A\) 在 \(F_i\) 上的投影。

在一般情形，采用与命题 12 的证明中相同的记号，我们有 \(\tau(\mathrm{E}_{i},\mathrm{F}_{i})=\tau(\mathrm{E}_{i},\mathrm{F}_{i}/\mathrm{N}_{i})\) 与 \(\tau(\mathrm{E},\mathrm{F})=\tau(\mathrm{E},\mathrm{F}/\mathrm{N})\)（IV，第 2 页），且由于 F/N 是诸 \(F_{i}/N_{i}\) 的拓扑直和，我们已归结为前面的情形。

证毕。

在本节余下部分，我们假定 \((\mathrm{E}_i)_{i\in \mathbf{I}}\) 是一个局部凸空间族。设 S 表示诸 \(\mathrm{E}_i\) 的拓扑直和，P 表示它们的乘积。我们定义一个线性映射 \(\theta :\mathrm{S}'\to \prod_{i\in \mathbf{I}}\mathrm{E}_i'\)，称为典范映射，其定义为

\[
\theta (x ^ {\prime}) = (x ^ {\prime} | \mathrm{E} _ {i}) _ {i \in \mathrm{I}} (x ^ {\prime} \in \mathrm{S} ^ {\prime})\tag{10}
\]

（其中 \(S'\) 表示 \(S\) 的对偶，\(E_i'\) 表示 \(E_i\) 的对偶）。

命题 14. --- (i) 映射 \(\theta\) 是从 \(S = \oplus E_i\) 的强（相应地，弱）对偶到诸 \(E_i\) 的强（相应地，弱）对偶的乘积上的一个同构：

\begin{enumerate}
\def\labelenumi{(\roman{enumi})}
\setcounter{enumi}{1}
\item
  为使 \(\mathbf{S}'\) 的子集 A 等度连续，必要且充分的是 \(\theta (\mathbf{A})\) 到 \(\mathrm{E}_i'\) 上的投影（对所有 \(i\in \mathbf{I}\)）等度连续。
\item
  Mackey 拓扑 \(\tau(\mathbf{S},\mathbf{S}')\) 是诸 Mackey 拓扑 \(\tau(\mathbf{E}_i,\mathbf{E}_i')\) 的直和拓扑。
\item
  拓扑 \(\beta (\mathbf{S},\mathbf{S}^{\prime})\) 是诸拓扑 \(\beta (\mathrm{E}_i,\mathrm{E}_i^{\prime})\) 的直和拓扑。
\end{enumerate}

\(\theta\) 的双射性立即由拓扑直和的定义得出（II，第 30 页，命题 6）。于是断言 (i) 对强拓扑由 IV 第 12 页的命题 12 得出，对弱拓扑由 II 第 50 页的命题 8 得出。类似地，(iii) 由（IV，第 12 页的）命题 13 得出，(iv) 由（IV，第 12 页的）命题 12 得出。

为证明 (ii)，设 A 是 S' 的一个子集。令

\[
q (x) = \sup _ {x ^ {\prime} \in \mathbf {A}} | \langle x, x ^ {\prime} \rangle | \quad \text { for } \quad x \in \mathrm{S};\tag{11}
\]

设 \(q_{i}\) 表示 q 在 \(E_{i}\) 上的限制，由此得

\[
q _ {i} (x _ {i}) = \sup _ {x _ {i} ^ {\prime} \in \mathbf {A} _ {i}} | \langle x _ {i}, x _ {i} ^ {\prime} \rangle | \quad \text { for } \quad x _ {i} \in \mathrm{E} _ {i},\tag{12}
\]

其中 \(A_{i}\) 表示 \(\theta(\mathbf{A})\) 在 \(E_{i}^{\prime}\) 上的投影。为使 A 等度连续，必要且充分的是 q 有限（即每个 \(q_{i}\) 有限）且连续。鉴于拓扑直和上连续半范数的刻画（II，第 27 页，命题 5），这等价于假定每个 \(q_{i}\) 连续，亦即每个集 \(A_{i}\) 等度连续。证毕。

设 \(\phi\) 是一个线性映射，称为典范映射，从 \(\oplus E_i'\) 到对偶 \(P'\)（

即 \(\mathbf{P} = \prod_{i\in \mathbf{I}}\mathrm{E}_i\) 的对偶）中，由公式

\[
\langle x, \phi (x ^ {\prime}) \rangle = \sum_ {i \in \mathbf {I}} \langle x _ {i}, x _ {i} ^ {\prime} \rangle\tag{13}
\]

（其中 \(x = (x_{i})\) 属于 \(\mathrm{P}\)，\(x^{\prime} = (x_{i}^{\prime})\) 属于 \(\bigoplus_{i\in \mathbf{I}}\mathrm{E}_{i}^{\prime}\)）定义。

命题 15. --- (i) 映射 \(\phi\) 是从诸 \(E_i\) 的强对偶的拓扑直和到 \(P = \prod_{i \in I} E_i\) 的强对偶上的一个同构。

\begin{enumerate}
\def\labelenumi{(\roman{enumi})}
\setcounter{enumi}{1}
\item
  为使 \(\mathbf{A}\)（\(\mathbf{P}'\) 的一个子集）等度连续，必要且充分的是它包含在一个有限和 \(\sum_{i\in \mathbf{J}}\phi (\mathbf{A}_i)\) 中，其中 \(\mathbf{J}\subset \mathbf{I}\) 有限，且 \(\mathbf{A}_i\) 在 \(\mathbf{E}_i'\) 中等度连续（对所有 \(i\in \mathbf{J}\)）。
\item
  Mackey 拓扑 \(\tau (\mathbf{P},\mathbf{P}^{\prime})\) 是诸拓扑 \(\tau (\mathrm{E}_i,\mathrm{E}_i')\) 的乘积拓扑。
\item
  拓扑 \(\beta (\mathbf{P},\mathbf{P}^{\prime})\) 是诸拓扑 \(\beta (\mathrm{E}_i,\mathrm{E}_i^{\prime})\) 的乘积拓扑。
\end{enumerate}

\(\phi\) 是单射这一点是直接的。\(\mathbf{P}\) 中 0 的一个基本邻域系由形如 \(\mathrm{V} = \prod_{i\in \mathbf{J}}\mathrm{V}_i\times \prod_{i\in \mathbf{I} - \mathbf{J}}\mathrm{E}_i\) 的集组成，其中 \(\mathrm{J}\subset \mathrm{I}\) 有限，且 \(\mathrm{V}_i\) 是 \(\mathrm{E}_i\) 中 0 的一个邻域（\(i\) 属于 \(\mathrm{J}\)）。\(\mathrm{V}\) 在 \(\mathbf{P}'\) 中的极集等于 \(\sum_{i\in \mathbf{J}}\phi (\mathrm{V}_i^\circ)\)。这就证明了 \(\phi\) 的满射性以及断言 (ii)。

断言 (i) 与 (iv) 由（IV，第 12 页的）命题 12 得出，(iii) 由（IV，第 12 页的）命题 13 得出。

推论. --- 桶型空间的任意乘积都是桶型的。

局部凸空间 E 是桶型的，当且仅当其初始拓扑与 \(\beta(\mathrm{E}, \mathrm{E}')\) 相同（IV，第 4 页，注 3）。因此只需应用命题 15 (iv)。

